% Independent audit and proof of the LM6/LM8 comparison requested in B.4.
% Attribution: AI-produced by GPT-6 Astra (OpenAI, ChatGPT Pro mode).

\subsection{The metabelian units and the extra arithmetic lifting condition}

\paragraph{LU.1 (Known, source opened): the precise metabelian input.}
Put $M=F/F''$, $A=M^{\rm ab}=\widehat{\mathbf Z}^2$,
$R=\widehat{\mathbf Z}[[A]]$, and $c=[x,y]$. The continuous
Alexander module is $M'=Rc$, free of rank one. Moreover
\[
 \operatorname{IOut}(M)=R^{\times,\epsilon=1}/A,
 \qquad
 \operatorname{Out}(M)=
 (R^{\times,\epsilon=1}/A)\rtimes\operatorname{GL}(A),
\]
with the natural action on the completed group ring. The
determinant of an IA automorphism is its multiplier on $Rc$;
an IA automorphism is inner exactly when this determinant is a
monomial in $A$. These are Chen--Deligne, Corollary 4.1.5,
Proposition 4.3.7 and Theorem 5.2.4. Their Corollary 4.1.6 also
gives $Z(M)=1$. The cited statements concern this actual
profinite quotient, rather than the completion of a discrete
formula with an unverified comparison map.

\paragraph{LU.2 (Proved): LM6 A.10 and A.49 are the units of LM8 A.22.}
There is a continuous injective homomorphism with closed image
\[
 \widehat{\mathbf Z}^{\times}\longrightarrow
 C_{\operatorname{Out}_{\rm lac}(M)}(G_K),\quad
 t\longmapsto[\phi_t]
\]
for every number field $K$. Its image is IA, centralizes the
geometric $S_3$, and is identity in every maximal pro-prime
quotient. The coefficient-five element $t_5=-1$, $t_p=1$ for
$p\ne5$, is exactly the involution of LM8 A.22. The analogous
choices for arbitrary sets of primes $p\ge5$ give LM6 A.10.

\emph{Proof.}
For $p\ne3$, let $A^{(p')}\to C_3$ send both named generators
to the same generator $t_0$, and let $H$ be its kernel.
Normalized averaging on every finite image of $H$ gives a
compatible idempotent $e_H\in\mathbf Z_p[[A^{(p')}]]$, because
all its averaging denominators are prime to $p$. Define
\[
 e_p=e_H\left(1-\frac{1+t_0+t_0^2}{3}\right)
\]
in the coefficient-$p$ factor, and zero in all other factors.
It is nonzero and has augmentation zero. On its summand,
the prime-to-$p$ parts of $X,Y,XY$ are nontrivial cubic roots;
therefore $X-1,Y-1,XY-1$ are units there. The kernel $H$ and
the pair of nontrivial characters are invariant under scalar
units and under the peripheral $S_3$. This construction applies
also at $p=2$.

For $p=3$ use $B=(\mathbf Z/4)^2$ and the six characters
\[
 \mathcal S=\{(a,b):a\ne0,\ b\ne0,\ a+b\ne0\pmod4\}.
\]
Their Fourier projector belongs to $\mathbf Z_3[B]$: conjugate
characters pair, and its coefficients have denominator dividing
eight. Average over the kernel of $A^{(3')}\to B$ to lift it
to an idempotent $e_3$ in the coefficient-three factor of $R$.
The set $\mathcal S$ is stable under $(a,b)\mapsto(b,a)$,
$(a,b)\mapsto(-a-b,b)$ and $(a,b)\mapsto(-a,-b)$; hence $e_3$
has the same invariance properties. Its augmentation is zero,
and $X-1,Y-1,XY-1$ are units on its summand.

For $t=(t_p)_p$ define, component by component,
\[
 u_t=1+\sum_p(t_p-1)e_p,\qquad
 r_t=\frac{u_t-1}{(X-1)(Y-1)},\qquad
 s_t=\frac{X(Y-1)r_t}{XY-1}.
\]
Each fraction is computed only on the indicated idempotent
summands and is extended by zero elsewhere. Put
$f_t=c^{r_t}$, $h_t=c^{s_t}$, where the exponents denote the
continuous $R$-module structure. The actual endomorphism
\[
 \phi_t(x)=x,\qquad \phi_t(y)=f_t^{-1}yf_t
\]
exists by the universal property of the free profinite
metabelian group. Direct multiplication gives
\[
 \phi_t(c)=c^{u_t},\qquad
 \phi_t(xy)=c^{X(Y-1)r_t}xy=h_t^{-1}xyh_t.
\]
It preserves all three labelled peripheral classes. Since it
is IA, it acts $R$-linearly on $M'$, and the conjugator for
$\phi_t\phi_v(y)$ has exponent
\[
 r_t+u_tr_v
   =\frac{u_tu_v-1}{(X-1)(Y-1)}=r_{tv}.
\]
Consequently $\phi_t\phi_v=\phi_{tv}$ and $\phi_{t^{-1}}$
is its inverse; these are actual automorphisms. In particular
$t_p=-1$ gives determinant $1-2e_p$ and an involution.

Arithmetic acts on $A$ as the scalar $\chi(\sigma)$. In LU.1,
conjugation by an arbitrary pair $(v,g)$ sends $(u,1)$ to
$(g(u),1)$, because the IA unit group is abelian. Invariance
of every $e_p$ thus proves outer commutation with $G_K$ and
with $S_3$. This step does not assert equality of arbitrarily
chosen representatives in $\operatorname{Aut}(M)$.

Projection to the coefficient-$p$ maximal pro-$p$ Alexander
module kills every $e_q$: for $q\ne p$ its coefficient is
killed, and for $q=p$ its nontrivial prime-to-$p$ character
becomes trivial. Thus the actual automorphism becomes the
identity in each maximal pro-prime quotient. If $u_t$ were a
monomial, these projections would make each pro-$p$ component
of that monomial trivial; the monomial would be one. The
nonzero torsion-free module $e_pR$ would then force $t_p=1$
for all $p$. LU.1 therefore proves outer injectivity. Every
finite group quotient uses finitely many coefficient primes
and residue digits, giving continuity. Compactness gives a
closed image.

The new independent program \texttt{check\_metabelian\_units.py}
checks the coefficient-three Fourier idempotent, its square,
augmentation and all stated symmetries exactly. It also constructs
the finite quotients of orders 48, 75 and 324 and checks generation,
every multiplication pair, all possible inner conjugators and the
three peripheral identities. These checks include both primes
excluded from the elementary cubic-character description at five.
They certify the displayed finite images; the inverse-limit
construction above proves the infinite assertion. \hfill$\square$

\paragraph{LU.3 (Proved): the center-free comparison needed one derived step higher.}
For every $m\ge2$, the group $S_m=F/F^{(m)}$ is center-free.
Hence its arithmetic extension over a number field $K$ is
canonically the pullback
\[
 E_m\simeq\operatorname{Aut}(S_m)
             \times_{\operatorname{Out}(S_m)}G_K.
\]
An outer automorphism centralizing $G_K$ therefore induces an
actual continuous automorphism of $E_m$ over $G_K$.

\emph{Proof.}
The base case is LU.1. Inductively a central element of $S_m$
lies in $F^{(m-1)}/F^{(m)}$. Let $H$ range over the open normal
subgroups of $F$ containing $F^{(m-1)}$. Its image $z_H$ in
$H^{\rm ab}$ is fixed by $F/H$ and maps to zero in $F^{\rm ab}$.
The transfer identity therefore gives
\[
 [F:H]z_H=\operatorname{Ver}_{F/H}(0)=0.
\]
Profinite Nielsen--Schreier makes $H^{\rm ab}$ a finite-rank
free profinite abelian group, which has no nonzero integer
torsion; hence $z_H=0$.
Moreover $\bigcap_HH'=F^{(m)}$: an element of $F^{(m-1)}$
outside $F^{(m)}$ survives in a finite quotient of $S_m$;
the inverse image of the abelian image of $F^{(m-1)}$ is one
such $H$, and its derived subgroup cannot contain the chosen
element. Thus the central element is trivial.

The map from $E_m$ to the displayed pullback sends an arithmetic
element to its conjugation action and its image in $G_K$.
Its kernel is $Z(S_m)=1$, and its image contains every fiber:
any two representatives of a prescribed outer action differ by
an inner automorphism. It is therefore a continuous isomorphism
of profinite groups. Conjugating the first component by a
representative of the commuting outer class defines the claimed
extension automorphism. \hfill$\square$

\paragraph{LU.4 (Proved): audit of LM6 A.56, including the constant twists.}
For every number field $K$, the image in $\operatorname{Out}(M)$
of the centralizer of $G_K$ in $\operatorname{Out}(S_3)$ contains
no nontrivial IA class. Here $S_3=F/F'''$ denotes the third
derived quotient, not the symmetric group. In particular every
nonidentity unit in LU.2 is obstructed one derived step higher
if commutation with $G_K$ is required.

\emph{Proof.}
Let $[\psi]$ be such a class and let $u\in R^{\times,\epsilon=1}$
represent its metabelian determinant. Write $I=\ker\epsilon$.
The continuous map $u\mapsto u-1\bmod\overline{I^2}$ identifies
$I/\overline{I^2}$ with $A$: $X^aY^b$ has linear term
$a(X-1)+b(Y-1)$, and the square-zero ring
$\widehat{\mathbf Z}\oplus A$ proves the converse independence.
Multiplying by the inverse monomial with this linear term
normalizes $u$ to a unit $u_0$ with zero linear term.

Outer commutation says $u_0/\sigma_R(u_0)\in A$ for every
$\sigma\in G_K$. This monomial has zero linear term and is
therefore one. Thus $\sigma_R(u_0)=u_0$. Inner adjustments lift
through $S_3\to M$, so choose an actual representative whose
metabelian determinant is $u_0$. By LU.3 it extends to an
automorphism $\beta$ of $E_3$ over $G_K$.

This choice makes $\beta$ the identity on $E_1$. Indeed if
$b_\sigma$ is conjugation by an arithmetic lift on $M$, its
action on $Rc$ is $r\mapsto\sigma_R(r)v_\sigma$. The commutator
of our representative with $b_\sigma$ acts on $Rc$ by
$u_0/\sigma_R(u_0)=1$. It is inner, say by $d_\sigma\in M$.
An inner action on $Rc$ multiplies by the monomial
$d_\sigma^{\rm ab}$, and its faithfulness gives $d_\sigma\in M'$.
The correction of each arithmetic lift thus vanishes in $E_1$;
the geometric abelian action was already the identity.

Suppose $u_0\notin A$. The closed procyclic subgroups
$C=\widehat{\mathbf Z}c$ and $D=\widehat{\mathbf Z}u_0c$ in
$M'$ cannot be conjugate. A conjugator acts by a monomial
$a\in A$, so conjugacy would give $u_0=\lambda a$ with
$\lambda\in\widehat{\mathbf Z}^{\times}$; augmentation forces
$\lambda=1$. Compactness detects their nonconjugacy in a
finite quotient: otherwise the nonempty closed sets of
conjugators modulo every open normal subgroup would have
the finite-intersection property and yield an actual conjugator.
Intersecting the finite automorphism-orbit of a detecting
open normal subgroup produces a stable finite metabelian
quotient $q:M\to G$ with induced automorphism $\gamma$ and
nonconjugate equal-order cyclic groups
$J=\langle q(c)\rangle$, $J'=\gamma(J)$. For $g=q(c)$,
\[
             \#(G/J)^g>0,\qquad \#(G/J')^g=0.
\]

After a finite extension of $K$, the geometric $G$-torsor
$Y\to U_0=\mathbf P^1\setminus\{0,1,\infty\}$ descends with
constant deck group and an arithmetic quotient $G\times G_K$.
On this quotient $\beta$ has form
$(h,\sigma)\mapsto(\gamma(h)z_\sigma,\sigma)$, with
$z:G_K\to Z(G)$ a continuous homomorphism. A further finite
extension kills $z$. We may also replace $K$ by its Galois
closure; this does not change the identity on $E_1$.

Hilbert irreducibility supplies $a\in U_0(K)$ whose fiber
$Y_a$ is a $G$-Galois field extension $L/K$; this is precisely
Serre, Proposition 3.3.1 and Theorem 3.4.1, in the opened
second edition. Lift $g$ to the normal closure of $L/\mathbf Q$.
Chebotarev, as stated in Milne, \emph{Class Field Theory},
Chapter V, Theorem 3.23, gives a rational prime $q$ split in
$K$ and a degree-one place $v$ whose fiber Frobenius is $g$.
Exclude the finitely many primes of bad marked reduction,
those dividing $|G|$, and those where $a$ or $1-a$ is not a
unit. Put $N=q-1$ and $K'=K(\mu_N)$. The place extends with
residue field $\mathbf F_q$, since $q\equiv1\pmod N$.

Use the actual constant-twisted abelian cover
\[
 B:\quad z_1^N=t/a,\qquad z_2^N=(1-t)/(1-a).
\]
Its arithmetic monodromy factors through $E_1$, so $\beta$
preserves it exactly, including its two Kummer constants.
Its only nonempty rational fiber over $\mathbf F_q$ is above
$a_0=a\bmod v$, where it has $N^2$ points. Therefore the
covers $Z=(Y/J)\times_{U_0}B$ and
$Z'=(Y/J')\times_{U_0}B$ satisfy
\[
 \#Z(\mathbf F_q)=N^2\#(G/J)^g>0,
 \qquad \#Z'(\mathbf F_q)=N^2\#(G/J')^g=0.
\]

On the other hand their geometric monodromy is metabelian.
Every component subgroup $H\subset F$ contains $F''$, hence
$F'''\subset H'$. Thus the quotient $F/F'''$ retains exactly
the full abelianization $H^{\rm ab}$ required for its first
cohomology. The automorphism $\beta$ identifies the arithmetic
cover data, the component permutation modules, and these
abelianizations, giving equivariant identifications of both
$H^0$ and $H^1$ of $Z_{\overline K}$ and $Z'_{\overline K}$.
For any coefficient prime different from $q$, good marked
reduction transports these representations to the special
fiber. The normalized products have good marked reduction:
locally at a boundary point their equations are tame Kummer
equations, and normalization replaces the indices by their
least common multiple, still prime to $q$. Poincare duality
gives equality on $H_c^1$; $H_c^0=0$ and $H_c^2$ is the
common component module twisted by $-1$. The trace formula
therefore forces equal point counts, a contradiction. The
cohomology facts used here are Milne, \emph{Lectures on Etale
Cohomology}, Theorems 24.1 and 29.4, together with proper
smooth base change for the smooth compactifications and their
finite etale boundaries.

It follows that $u_0\in A$, and LU.1 makes the metabelian IA
outer class trivial. Notice the exact scope: the argument
does not prove that the kernel of
$C_{\operatorname{Out}(F/F''')}(G_K)\to\operatorname{Out}(M)$
is trivial. \hfill$\square$

\paragraph{LU.5 (Proved): audit of LM6 A.57--A.58.}
The abelian image of the peripheral metabelian centralizer
of $G_K$ is exactly the geometric symmetric group. Consequently
\[
 \operatorname{image}\left(
 C_{\operatorname{Out}_{\rm lac}(F/F''')}(G_K)
       \longrightarrow\operatorname{Out}(M)\right)
       =S_3^{\rm geom},
\]
and the same holds for every higher derived length projecting
to $M$.

\emph{Proof.}
An abelian automorphism preserving the three primitive
peripheral lines is a scalar unit times a geometric permutation:
fixing the first two lines makes it diagonal and fixing their
sum makes the diagonal entries equal. Divide by the geometric
permutation and write its common multiplier $\rho$.

We need an exact arithmetic normalization. In $M$ the normalizer
of $\langle x\rangle$ is $\langle x\rangle$. Indeed a centralizer
element has zero $y$-coordinate by the free class-two quotient,
and hence is $x^a c^r$. Commutation says $(X-1)r=0$. In every
coefficient-$p$ factor, refining the first cyclic group coordinate
by $p^j$ makes an $X$-invariant coefficient project to $p^j$
equal summands; as $j$ is arbitrary it is zero. This proves
injectivity of $X-1$ on the full completed ring. Normalizing
$\langle x\rangle$ already forces exponent one by abelianization,
so the normalizer equals the centralizer. The same holds for $y$.

The tangential sections at $0$ towards $1$ and at $1$ towards
$0$ have the same image in $E_1=A\rtimes G_K$: on the covers
$t^{1/n}$ and $(1-t)^{1/n}$ their leading coefficients or values
are one, and choose the compatible roots one. Their arithmetic
normalizers are respectively $\langle x\rangle\rtimes G_K$
and $\langle y\rangle\rtimes G_K$. Write an extension
automorphism obtained from LU.3 as
$(v,g)\mapsto(\rho v+b(g),g)$ on $E_1$. Peripheral preservation
and the two normalizer equalities imply
$b_2(g)=(1-\chi(g))(a_0)_2$ and
$b_1(g)=(1-\chi(g))(a_1)_1$, where $a_i$ are the abelian
coordinates of its two geometric cusp conjugators. Thus $b$
is an actual coboundary. An inner adjustment gives exactly
$(v,g)\mapsto(\rho v,g)$.

Let $q$ be an odd rational prime split in a Galois closure of
$K$, $N=q-1$, and pass to $K'=K(\mu_N)$ at a degree-one place
over $q$. For each integer $a\in\{2,\ldots,q-1\}$ use
\[
 B_a:\ z_1^N=t/a,\ z_2^N=(1-t)/(1-a).
\]
If $r=\rho\bmod N$, the exact abelian action sends its
arithmetic cover subgroup to that of
\[
 B'_a:\ w_1^N=t/a^r,\ w_2^N=(1-t)/(1-a)^r.
\]
This follows by applying $(v,g)\mapsto(\rho v,g)$ to the
equations $v=(\kappa_N(a)(g),\kappa_N(1-a)(g))$ defining the
cover subgroup. Their geometric monodromy is abelian, so
$F''$ lies in each component subgroup's commutator and already
$E_2$ retains its full first cohomology. The same cohomology
and trace argument as LU.4 forces equal point counts. Both
covers have good marked reduction, being Fermat curves of
degree $q-1$ with unit coefficients and their finite etale
boundaries removed. The first has $N^2$ rational points;
the second has $N^2$ precisely when
\[
                       a^r+(1-a)^r=1\quad\hbox{in }\mathbf F_q.
\]
This identity therefore holds for all $a\ne0,1$. The map
$f(0)=0$, $f(x)=x^r$ on $\mathbf F_q^\times$ is multiplicative
and bijective. Taking $a=x/(x+y)$ proves additivity for nonzero
$x,y,x+y$. The remaining cases follow from $f(-1)=-1$, since
$r$ is odd. Thus $f$ is an automorphism of the prime field,
and is the identity. Hence $\rho\equiv1\pmod{q-1}$.
For every $n$, Chebotarev permits such a $q$ also split in
$\mathbf Q(\mu_n)$, so $\rho\equiv1\pmod n$; thus $\rho=1$.

The projected centralizer at length three consequently has
the abelian action of one geometric permutation. Dividing by
it leaves an IA class, trivial by LU.4. The geometric six
permutations are defined over $\mathbf Q$ and lift at every
derived level. This proves equality, including higher levels
by projection to length three. \hfill$\square$

\paragraph{LU.6 (Proved): reconciliation with the three weight-five directions.}
LB.1 concerns arbitrary multiplier-one peripheral automorphisms
of the full free profinite group; it imposes no commutation
with arithmetic. LU.2 concerns actual metabelian outer classes
which commute with arithmetic, and LU.4 excludes lifting these
classes while retaining that extra condition. Their determinants
lie on nontrivial prime-to-$p$ character summands and have zero
image in every absolute pro-$p$ quotient, so they are not any
of the nonzero weight-five directions. Thus there is neither
a conflict of categories nor a counterexample to R1. R1 retains
three peripheral conjugators, but its hypotheses do not include
Galois commutation. At primes two and three LU.2 and LU.4 apply
unchanged; the invariant pro-$\ell$ lifting step of LP.2 is
only asserted for $\ell\ge5$, exactly where division by six
and the two coprime complement arguments are valid.
